% Generated draft: requires mathematical and source review.
\documentclass[11pt]{article}
\usepackage[margin=25mm]{geometry}
\usepackage{fontspec}
\setmainfont{Latin Modern Roman}
\usepackage{amsmath,amssymb,mathtools,unicode-math}
\setmathfont{Latin Modern Math}
\usepackage{enumitem,needspace,xurl,hyperref}
\setlength{\parindent}{0pt}
\setlength{\parskip}{5pt}
\setlength{\emergencystretch}{3em}
\newcommand{\N}{\mathbb{N}}
\newcommand{\Z}{\mathbb{Z}}
\newcommand{\Q}{\mathbb{Q}}
\newcommand{\R}{\mathbb{R}}
\newcommand{\C}{\mathbb{C}}
\newcommand{\F}{\mathbb{F}}
\newcommand{\A}{\mathbb{A}}
\newcommand{\PP}{\mathbb{P}}
\newcommand{\E}{\mathbb{E}}
\newcommand{\eps}{\varepsilon}
\newcommand{\dd}{\,\mathrm{d}}
\begin{document}
\section*{UnsolvedMath ID: 701049}
\textbf{Source problem code:} catalogue entry
\subsection{Definitions and mathematical statement}
For a group \(G\), define its lower central series by \(\gamma_1(G)=G\) and
\[
\gamma_{i+1}(G)=[\gamma_i(G),G]\qquad(i\geq 1),
\]
where \([H,G]\) is the subgroup generated by \(h^{-1}g^{-1}hg\) for \(h\in H\) and \(g\in G\). A nontrivial nilpotent group has nilpotency class \(c\) if \(\gamma_c(G)\neq\{1\}\) and \(\gamma_{c+1}(G)=\{1\}\). Using the conventional definition, a finite \(5\)-group \(G\) of order \(5^n\), with \(n\geq 2\), has maximal class when its nilpotency class is \(n-1\).

\textbf{Classification problem (M. F. Newman, as attributed in the statement).} Classify, up to group isomorphism, all finite \(5\)-groups of maximal class. Equivalently, describe all isomorphism classes of groups \(G\) satisfying
\[
|G|=5^n,\qquad \gamma_{n-1}(G)\neq\{1\},\qquad \gamma_n(G)=\{1\},\qquad n\geq 2.
\]
Here classification is interpreted as classification up to isomorphism. The stated convention includes groups of order \(25\); the source does not specify a lower-order cutoff.

The source says that computer calculations suggest some conjectures and refers to Newman's 1990 contribution, but it does not state those conjectures. They cannot be formulated separately from the supplied record. Newman's supplied comment of 2025 points to related work by Cant, Dietrich, Eick, and Moede (2022), without asserting that this classification problem is resolved. The registry labels the problem open; this is a supplied status, not an independently verified assessment. The registry proposer attribution to A. Moret\'{o} conflicts with the attribution to M. F. Newman in the statement.
\subsection{Short English statement}
Describe every finite group of order \(5^n\) and nilpotency class \(n-1\), up to isomorphism, using the convention \(n\geq 2\). The source mentions computational conjectures and related work without specifying those conjectures or claiming a complete classification. Its proposer attributions conflict.
\subsection{Sources}
\begin{itemize}
\item Local UnsolvedMath registry, permanent ID \(701049\), source problem code \texttt{catalogue entry}, title \emph{Kourovka Notebook Problem 14.64}. The supplied record labels the problem open and lists A. Moret\'{o} as proposer, with proposed year 1999; the statement instead attributes it to M. F. Newman.
\item \emph{The Kourovka Notebook (Unsolved Problems in Group Theory)}, problem 14.64, 14th issue (1999). The supplied background identifies the 21st issue edition (2026), arXiv identifier \texttt{1401.0300v48}. The supplied collection description names E. I. Khukhro and V. D. Mazurov as editors and describes coverage of Issues 1--21 (1965--2026). These bibliographic identifications have not been independently verified.
\item M. F. Newman, contribution in \emph{Groups--Canberra, 1989}, Lecture Notes in Mathematics, vol. 1456, Springer, Berlin, 1990, pp. 49--62. The supplied statement cites this contribution in connection with conjectures suggested by computer calculations; no contribution title or explicit conjectures are supplied.
\item A. Cant, H. Dietrich, B. Eick, and T. Moede, \emph{Journal of Algebra}, vol. 604 (2022), pp. 429--450. M. F. Newman's supplied comment of 2025 describes this as work on related conjectures. No article title or more specific relationship to the classification target is supplied.
\end{itemize}
\end{document}
